% !TEX program = pdflatex
\documentclass[12pt]{article}

\usepackage[utf8]{inputenc}
\usepackage{CJKutf8}
\usepackage{amsmath,amssymb}
\usepackage{amsthm}
\usepackage{geometry}
\geometry{
  a4paper,
  top=25mm,
  bottom=25mm,
  left=20mm,
  right=20mm
}
\usepackage{xcolor}
\usepackage{url}
\usepackage[unicode,colorlinks=true,linkcolor=blue,urlcolor=blue]{hyperref}
\usepackage{xurl}
\Urlmuskip=0mu plus 2mu
\sloppy
\usepackage{tocloft}

% 目录中 section 条目使用点线填充到页码
\renewcommand{\cftsecleader}{\cftdotfill{\cftdotsep}}

% 基本定理环境（工作笔记：形式系统风格）
\newtheorem{definition}{定义}[section]
\newtheorem{assumption}{假设}[section]
\newtheorem{theorem}{定理}[section]
\newtheorem{proposition}{命题}[section]
\newtheorem{corollary}{推论}[section]
\newtheorem{lemma}{引理}[section]
\newtheorem{remark}{备注}[section]
\newtheorem{Certificate}{证书}[section]

\newcommand{\Jac}{\operatorname{Jac}}
\newcommand{\dd}{\mathrm{d}}
\newcommand{\code}[1]{\path{#1}}
\newcommand{\GFramework}{\textsc{G-Framework}}

\newcommand{\Ob}{\operatorname{Ob}}
\newcommand{\ObTot}{\operatorname{ObTot}}
\newcommand{\Hole}{\operatorname{Hole}}
\newcommand{\Lift}{\operatorname{Lift}}
\newcommand{\Cert}{\operatorname{Cert}}
\newcommand{\Loss}{\mathcal{L}}
\newcommand{\Flow}{\operatorname{Flow}}
\newcommand{\Upd}{\operatorname{Upd}}

% 避免少量不可控的换行溢出（尤其是混排 CJK 与等宽字体时）
\setlength{\emergencystretch}{6em}

% 对话稿默认取消首行缩进，并留出段间距，提升可读性
\setlength{\parindent}{0pt}
\setlength{\parskip}{0.6em}

\begin{document}
\begin{CJK*}{UTF8}{gkai}

\title{纤维--截面--孔洞--障碍量--势函数--条件化收敛\\
（独立说明稿：与层 $l_n$ 同伦塔链的对齐与收敛模板，形式逻辑版）}
\newcommand{\CertificateAuthorCN}{Gao Zheng（高政）}
\author{\CertificateAuthorCN}
\date{2026-03-13}
\maketitle

\section*{前言}
\addcontentsline{toc}{section}{前言}

本文的目标不是把纤维、截面、孔洞、障碍量、势函数与条件收敛分别作为孤立术语
处理，而是在 \GFramework{} 的障碍修复链条中，把它们组织成一个从层结构到收敛
判据的形式系统。若这些对象分散出现，孔洞容易只被看成拓扑缺口，势函数容易只
被看成数值代价，条件收敛容易只被看成附加假设；本文将三者统一为“障碍如何被
度量、如何被扭曲、如何在条件下收敛”的问题。

本文的第一层立场是层、纤维与截面的对象化。系统不是单一平面上的点集，而是由
层结构、遗忘投影、纤维与截面共同组成的多层对象。截面给出可选择的局部代表，
纤维给出被投影隐藏的内部自由度，遗忘投影则暴露出何处可能产生孔洞、崩溃与
上同调障碍。由此，障碍不再只是外部失败，而是层间投影与截面选择不相容时产生的
结构性失闭合。

本文的第二层立场是把孔洞转化为可度量的障碍量。孔洞的意义不只在于“缺了一块”，
而在于它阻断了截面延拓、破坏了同伦闭合，并诱导可被势函数读取的障碍量。势函数
不是任意打分，而是将孔洞、障碍类、同伦扭曲代价与修复进展压入同一个数值结构，
使“是否可修复”与“如何趋近修复”能够被比较。

本文的第三层立场是截面失败的可定位化。一个截面无法全局延拓时，失败并不总是
来自整个空间的不可处理性，而可能来自某个孔洞、某个纤维缺口、某个层间不兼容或
某个被遗忘投影隐藏的自由度。本文将这些失败转写为可定位的障碍对象，再通过障碍量
和势函数将其纳入可比较的修复过程。这样，修复不再是对全局空间的盲目重构，而是
对“何处失闭合、何处需补纤维、何处需扭曲”的精确响应。

本文的第四层立场是势函数的桥接作用。势函数在本文中同时连接拓扑信息与动力信息：
一方面，它读取孔洞、上同调障碍和截面不连续；另一方面，它给出同伦扭曲、收敛
下降与条件化控制的数值接口。由此，势函数不是替代拓扑证明的打分器，而是把拓扑
障碍送入可计算修复过程的桥。

本文的第五层立场是条件化收敛。本文并不把所有障碍都无条件消解为收敛结论，而是
明确区分混合指标、同伦塔、势函数下降与条件化收敛之间的证书关系。只有当障碍量
被合适的扭曲、势函数下降与层间相容条件共同控制时，收敛才是可证明的；否则，
本文只保留边界而不把它写成定理结论。

本文的第六层立场是两种无穷的分离。层数无穷、修复塔无穷、截面延拓过程无穷，与
势函数下降的极限行为不是同一件事。本文通过混合指标与条件化收敛，把结构无穷、
过程无穷和数值极限分开处理。这样可以避免把“层级可继续展开”误写成“势能必然
收敛”，也避免把“势能下降”误写成“拓扑障碍必然消失”。

本文的第七层立场是为后续同伦命中提供低残差雏形。后续统计力学同伦命中会把目标
集合写成低残差集合，而本文已经给出较早的版本：孔洞、障碍量和势函数共同定义一个
可被压低的目标区域。也就是说，后续的“命中低残差集合”在这里先表现为“使障碍量与
势函数进入可接受阈值”的条件化收敛问题。

因此，本文在整体 \GFramework{} 中承担的是“孔洞--障碍量--势函数--条件收敛”的
中间锚定作用。它向前连接纤维丛、截面与上同调障碍，向后连接统计力学命中中的
低残差集合与可控隧穿路径。概括地说，本文把“孔洞是否阻断修复”这一问题，推进为
“障碍量能否经由势函数与同伦扭曲被条件化压低”的形式系统。

更具体地说，本文将“孔洞”从视觉化直觉转化为纤维--截面系统中的失败位置。孔洞
不是单纯缺口，而是截面无法穿过、投影无法恢复、局部代表无法拼合的地方。障碍量
则把这种失败量化，使后续修复不再只依赖拓扑存在性，而能进入势函数下降和阈值判定。

从方法论上看，本文是从同伦修复走向统计命中的过渡稿。前者强调障碍类如何被吸收，
后者强调目标集合如何被命中；本文则把障碍写成势函数和条件化收敛对象，使“修复”
与“命中”之间出现可计算中介。后续低残差集合的概念，在这里已经以低障碍量、
低势能区域的形式出现。

从整体谱系看，本文避免两种过度说法：其一，孔洞存在并不等于不可修复；其二，势函数
下降也不等于拓扑障碍自动消失。只有当纤维、截面、障碍量、同伦扭曲和条件收敛证书
共同闭合时，本文才给出收敛结论。这种谨慎边界使它成为 \GFramework{} 修复链中的
重要中层控制稿。

\setcounter{tocdepth}{2}
\setcounter{secnumdepth}{2}
\tableofcontents
\clearpage

% ------------------------------
% 本文为独立说明稿：不依赖专著正文的外部引用。
% ------------------------------

\section{目标、写作约定与总体结构}

\begin{remark}[形式逻辑写作约定]
本文统一采用“形式逻辑 + 证书”体例，并执行以下规范：
\begin{itemize}
  \item \textbf{公理降级：}凡需要作为基础约束的强陈述，统一写成“公理（降级为定理）”，并配套\textbf{数学归纳法证书}与证明；
  \item \textbf{定理/引理/命题/推论：}每条编号断言均配套“证书 + 证明（基于数学符号推演逻辑的谓词因果链）”；
  \item \textbf{备注：}仅承担动机、边界与解释说明，不承担证明义务；
  \item \textbf{章节编号：}全部依赖 LaTeX 自动编号，不手工维护章节序号。
\end{itemize}
\end{remark}

本文把两条叙述链严格对齐，并给出可审计的“条件化收敛”证明骨架。

\begin{itemize}
\item \textbf{Chain-A（骨架链）}：
\[
\text{纤维}\to\text{截面}\to\text{孔洞}\to\text{障碍量}\to\text{势函数}\to\text{条件化收敛}.
\]
\item \textbf{Chain-B（层化链）}：
\[
\begin{aligned}
l_2
&\to(\text{连续统（一离散）})\to(\text{崩溃/上同调障碍})\\
&\to(\text{同伦塔度量，}l_{\ge3})\to(\text{同伦扭曲})\\
&\to(\text{变分法收敛})\to(\text{离散统（一连续）}).
\end{aligned}
\]
\end{itemize}

\begin{remark}[强结论边界]
本文所有“\textbf{强结论}”（例如“全局最优路径收敛”或“孔洞 $\Leftrightarrow$ 单一障碍类非零”）若缺乏对应假设，将被明确标注为\textbf{未证明}。
\end{remark}

\section{层、遗忘投影、纤维与截面}

\subsection{层与状态空间}

\begin{definition}[层与状态空间]\label{def:layer-space}
对每个整数 $n\ge 2$，给定一层 $l_n$ 及其可表述对象（系统状态）空间 $\mathcal S_n$。
层 $l_2$ 仅保留外延（extensional）信息；层 $l_{n\ge3}$ 允许加入语义/同伦/曲率/生成性等增强信息。
\end{definition}

\subsection{遗忘投影与纤维冗余}

\begin{definition}[遗忘投影]\label{def:forget-map}
对任意 $n\ge3$，给定遗忘投影
\[
\pi_{n\to2}:\mathcal S_n\to\mathcal S_2,
\]
其含义是把 $l_n$ 层对象遗忘到 $l_2$ 层可见的外延数据。
\end{definition}

\begin{definition}[纤维（冗余集合）]\label{def:fiber}
对任意 $x\in\mathcal S_2$，定义 $x$ 在第 $n$ 层的纤维
\[
\operatorname{Fib}_n(x):=\pi_{n\to2}^{-1}(x)\subseteq \mathcal S_n.
\]
当 $|\operatorname{Fib}_n(x)|>1$ 时，称 $l_2$ 对应于 $x$ 的信息对 $l_n$ 是\emph{冗余的}：即多个语义对象在遗忘后不可区分。
\end{definition}

\begin{remark}
“连续统（一离散）”可被解释为：在 $l_2$ 层输出的是离散外延标签（例如某个离散不变量），而其纤维 $\operatorname{Fib}_n(x)$ 可能含有不可数多的语义提升候选。
\end{remark}

\subsection{可接受集、可提升性与截面}

\begin{definition}[可接受集]\label{def:acceptable}
对每个 $n\ge3$，给定可接受集（满足一致性/约束/可审计性）
\[
\mathcal A_n\subseteq \mathcal S_n.
\]
\end{definition}

\begin{definition}[可接受提升集合]\label{def:lift}
对任意 $x\in\mathcal S_2$，定义
\[
\Lift_n(x):=\mathcal A_n\cap \operatorname{Fib}_n(x).
\]
若 $\Lift_n(x)\neq\varnothing$，则称 $x$ 在层 $l_n$ 上\emph{可提升}（存在满足约束的语义代表）。
\end{definition}

\begin{definition}[截面（冗余剔除的选择）]\label{def:section}
设
\[
\mathcal D_n:=\{x\in\mathcal S_2:\Lift_n(x)\neq\varnothing\}.
\]
称映射
\[
\sigma_n:\mathcal D_n\to \mathcal A_n
\]
为第 $n$ 层的截面，若
\[
\pi_{n\to2}\circ \sigma_n=\mathrm{id}_{\mathcal D_n}.
\]
\end{definition}

\begin{proposition}[截面存在与可提升选择规则等价]\label{prop:section-choice}
对固定 $n\ge 3$，存在截面 $\sigma_n$ 当且仅当对每个 $x\in\mathcal D_n$ 能从 $\Lift_n(x)$ 中选取一个元素作为代表。
\end{proposition}

\begin{Certificate}[证书：选择映射与恒等复合判据]\label{cert:section-choice}
证书包含两项：
\begin{enumerate}
\item 选择函数条款：$\forall x\in\mathcal D_n,\ \exists y_x\in\Lift_n(x)$；
\item 复合校验条款：$\pi_{n\to2}(y_x)=x$。
\end{enumerate}
目标是验证“可选择”与“存在截面”互相推出。
\end{Certificate}

\begin{proof}[证明（谓词因果链）]
由证书~\ref{cert:section-choice}：若存在选择规则 $x\mapsto y_x\in\Lift_n(x)$，定义
\[
\sigma_n(x):=y_x.
\]
由于 $y_x\in\Lift_n(x)\subseteq\pi_{n\to2}^{-1}(x)$，故
\[
\pi_{n\to2}(\sigma_n(x))=\pi_{n\to2}(y_x)=x,
\]
即 $\pi_{n\to2}\circ\sigma_n=\mathrm{id}_{\mathcal D_n}$，截面存在。
反向，若截面存在，则取 $y_x:=\sigma_n(x)$，显然 $y_x\in\Lift_n(x)$，给出选择规则。
故两者等价。
\end{proof}

\begin{remark}
截面 $\sigma_n$ 的存在把“纤维多值冗余”压缩为“可审计单值代表”。这正是 Chain-A 中“纤维 $\to$ 截面”的接口含义。
\end{remark}

\section{孔洞、崩溃与上同调障碍}

\subsection{孔洞判据（可行域为空）}

\begin{definition}[孔洞]\label{def:hole}
对任意 $x\in\mathcal S_2$，称 $x$ 是（第 $n$ 层）孔洞点，若
\[
\Lift_n(x)=\varnothing.
\]
记该性质为 $\Hole_n(x)$。
\end{definition}

\begin{remark}
“崩溃（上同调障碍）”在最小可审计意义下可取为：$\Lift_n(x)=\varnothing$（可行域为空）。此时“崩溃”不依赖修辞理解，而是一个可计算判据。
\end{remark}

\subsection{障碍类与完备性条款}

\begin{definition}[塔式障碍类与总障碍]\label{def:obstruction-tower}
设存在与提升塔相配套的一族障碍类
\[
\operatorname{ob}_{k+1}(x)\in H^{k+1}(\mathcal B_x;\mathcal G_{k,x}),\qquad k\ge2,
\]
并定义总障碍向量
\[
\ObTot(x):=\bigl(\operatorname{ob}_3(x),\operatorname{ob}_4(x),\ldots\bigr)
\in \prod_{k\ge2} H^{k+1}(\mathcal B_x;\mathcal G_{k,x}).
\]
\end{definition}

\begin{remark}[参考：障碍理论]
上述“逐级延拓失败对应上同调障碍类”的形式化口径属于标准障碍理论框架，可参见 \cite{Hatcher2002,Whitehead1978}。
\end{remark}

\begin{assumption}[障碍理论完备性 $\mathbf{(OT\text{-}Complete)}$]\label{ass:ot-complete}
对任意 $x\in\mathcal S_2$，有充要式
\[
\Lift_\infty(x)\neq\varnothing
\iff
\forall k\ge2,\ \operatorname{ob}_{k+1}(x)=0,
\]
其中 $\Lift_\infty$ 表示在塔极限意义下的可接受提升集合。
\end{assumption}

\begin{lemma}[必要性：可提升 $\Rightarrow$ 障碍全为零]\label{lem:ob-necessary}
若 $\Lift_\infty(x)\neq\varnothing$，则 $\forall k\ge2,\ \operatorname{ob}_{k+1}(x)=0$。
\end{lemma}

\begin{Certificate}[证书：失败测度归零判据]\label{cert:ob-necessary}
证书登记规则：每一级障碍类 $\operatorname{ob}_{k+1}(x)$ 都刻画“从第 $k$ 级向第 $(k+1)$ 级延拓失败”的等价类；若延拓数据已存在，则对应失败类必须为零类。
\end{Certificate}

\begin{proof}[证明（谓词因果链）]
由证书~\ref{cert:ob-necessary}：
\[
\Lift_\infty(x)\neq\varnothing
\Rightarrow
\text{各级延拓均可实现}
\Rightarrow
\text{各级失败测度消失}
\Rightarrow
\forall k\ge2,\ \operatorname{ob}_{k+1}(x)=0.
\]
引理成立。
\end{proof}

\begin{lemma}[充分性需要完备性]\label{lem:ob-sufficient}
若 $\forall k\ge2,\ \operatorname{ob}_{k+1}(x)=0$，则在假设~\ref{ass:ot-complete} 下有 $\Lift_\infty(x)\neq\varnothing$。
\end{lemma}

\begin{Certificate}[证书：完备性条款调用]\label{cert:ob-sufficient}
证书仅调用假设~\ref{ass:ot-complete} 的“$\Leftarrow$”方向：
\[
\forall k\ge2,\ \operatorname{ob}_{k+1}(x)=0
\Rightarrow
\Lift_\infty(x)\neq\varnothing.
\]
\end{Certificate}

\begin{proof}[证明（谓词因果链）]
由证书~\ref{cert:ob-sufficient} 直接得到结论，引理成立。
\end{proof}

\begin{theorem}[孔洞与总障碍的充要对应]\label{thm:hole-iff-obtot}
在假设~\ref{ass:ot-complete} 下，
\[
\Hole_\infty(x)\iff \ObTot(x)\neq 0,
\qquad
\Lift_\infty(x)\neq\varnothing \iff \ObTot(x)=0.
\]
\end{theorem}

\begin{Certificate}[证书：必要性 + 充分性双桥接]\label{cert:hole-iff-obtot}
证书包含两条桥接：
\begin{enumerate}
\item 引理~\ref{lem:ob-necessary}：$\Lift_\infty\neq\varnothing\Rightarrow\ObTot=0$；
\item 引理~\ref{lem:ob-sufficient}：$\ObTot=0\Rightarrow\Lift_\infty\neq\varnothing$。
\end{enumerate}
再对可提升命题取逻辑否定得到孔洞命题。
\end{Certificate}

\begin{proof}[证明（谓词因果链）]
由证书~\ref{cert:hole-iff-obtot}：
\[
\Lift_\infty(x)\neq\varnothing \Longleftrightarrow \ObTot(x)=0.
\]
对该等价式取否定并使用
\(
\Hole_\infty(x)\Longleftrightarrow \Lift_\infty(x)=\varnothing
\)
得到
\(
\Hole_\infty(x)\Longleftrightarrow \ObTot(x)\neq 0
\)。
定理成立。
\end{proof}

\begin{corollary}[非孔洞点的可提升判据]\label{cor:non-hole-lift}
在假设~\ref{ass:ot-complete} 下，
\[
\neg\Hole_\infty(x)\iff \ObTot(x)=0.
\]
\end{corollary}

\begin{Certificate}[证书：定理~\ref{thm:hole-iff-obtot} 的逻辑重写]\label{cert:non-hole-lift}
证书只需把定理~\ref{thm:hole-iff-obtot} 的第一条等价式两侧同时取否定。
\end{Certificate}

\begin{proof}[证明（谓词因果链）]
由证书~\ref{cert:non-hole-lift}：
\[
\Hole_\infty(x)\iff \ObTot(x)\neq 0
\Rightarrow
\neg\Hole_\infty(x)\iff \ObTot(x)=0.
\]
推论成立。
\end{proof}

\begin{remark}
若不引入假设~\ref{ass:ot-complete}，则一般只能保证必要性（引理~\ref{lem:ob-necessary}），而充分性属于\textbf{未证明}。
\end{remark}

\section{同伦塔、混合指标与两种无穷}

\subsection{层阶与语义参数的混合指标}

\begin{definition}[混合指标]\label{def:mixed-index}
对每个 $n\ge3$，令 $\Theta_n$ 为语义/算子自由度参数空间，并假设
\[
|\Theta_n|\ge|\mathbb R|.
\]
则一次“语义增强步骤”的指标为
\[
(n,\theta_n)\in \mathbb N_{\ge3}\times\Theta_n.
\]
其中 $n$ 是可数层阶，而 $\theta_n$ 是不可数连续自由度。
\end{definition}

\begin{lemma}[两种无穷在同一指标集内共存]\label{lem:two-infinities}
固定任意 $\theta\in\Theta_n$ 时，
\(
\{(n,\theta):n\in\mathbb N_{\ge3}\}
\)
给出可数子结构；固定任意 $n$ 时，
\(
\{(n,\theta):\theta\in\Theta_n\}
\)
给出不可数子结构。
\end{lemma}

\begin{Certificate}[证书：两类单射嵌入]\label{cert:two-infinities}
证书登记两条单射：
\[
\iota_1:\mathbb N_{\ge3}\hookrightarrow \mathbb N_{\ge3}\times\Theta_n,\ n\mapsto(n,\theta),
\]
\[
\iota_2:\Theta_n\hookrightarrow \mathbb N_{\ge3}\times\Theta_n,\ \theta\mapsto(n,\theta).
\]
由此分别保留可数性与不可数性。
\end{Certificate}

\begin{proof}[证明（谓词因果链）]
由证书~\ref{cert:two-infinities}：
\(
\iota_1
\)
把可数集嵌入指标集，故存在可数子结构；
\(
\iota_2
\)
把不可数集嵌入指标集，故存在不可数子结构。
两种无穷在同一指标框架中共存。
\end{proof}

\subsection{算子与语义的唯一索引}

\begin{definition}[算子库与语义映射]\label{def:op-sem}
令 $\mathcal O_n$ 为第 $n$ 层边缘/修复算子库。给定语义映射
\[
\operatorname{Sem}_n:\mathcal O_n\to \Sigma_n,
\]
其中 $\Sigma_n$ 为语义对象（或语义类型）空间。
\end{definition}

\begin{assumption}[算子定义语义（单射性）]\label{ass:sem-inj}
假设 $\operatorname{Sem}_n$ 单射：若 $\operatorname{Sem}_n(o)=\operatorname{Sem}_n(o')$，则 $o=o'$。
\end{assumption}

\begin{proposition}[单射条件下的算子语义唯一索引]\label{prop:sem-unique-index}
在假设~\ref{ass:sem-inj} 下，每个语义值最多对应一个算子，即
\[
\forall s\in\Sigma_n,\ \bigl|\operatorname{Sem}_n^{-1}(\{s\})\bigr|\le 1.
\]
\end{proposition}

\begin{Certificate}[证书：单射定义的逆像基数判据]\label{cert:sem-unique-index}
证书目标：把“单射”改写为“任一像元的逆像基数不超过 1”，即
\[
\operatorname{Sem}_n\ \text{单射}
\Longleftrightarrow
\forall s,\ |\operatorname{Sem}_n^{-1}(\{s\})|\le 1.
\]
\end{Certificate}

\begin{proof}[证明（谓词因果链）]
由证书~\ref{cert:sem-unique-index}：
若存在 $s$ 使逆像中有 $o\neq o'$，则
\(
\operatorname{Sem}_n(o)=\operatorname{Sem}_n(o')=s
\)
与单射矛盾。
故每个语义值最多对应一个算子，命题成立。
\end{proof}

\begin{remark}
假设~\ref{ass:sem-inj} 的弱化版本是“到同余类的双射”：把语义等价与算子等价同时按同一同余关系取商。若不作任何假设，则“算子 $\equiv$ 语义”为\textbf{未证明}。
\end{remark}

\section{障碍量、同伦扭曲与势函数}

\subsection{同伦塔度量与障碍量}

\begin{definition}[度量与障碍量]\label{def:ob-measure}
对每个 $n\ge3$，在 $\mathcal S_n$ 上给定度量（或伪度量）$d_n$。
定义障碍量为到可接受集的距离：
\[
\Ob_n(y):=\inf_{z\in\mathcal A_n} d_n(y,z)\in\mathbb R_{\ge0}.
\]
\end{definition}

\begin{proposition}[闭集情形下的零障碍判据]\label{prop:ob-zero-criterion}
若 $\mathcal A_n$ 在 $d_n$ 诱导拓扑下为闭集，则
\[
\Ob_n(y)=0\iff y\in\mathcal A_n.
\]
\end{proposition}

\begin{Certificate}[证书：距离到闭集的标准判据]\label{cert:ob-zero-criterion}
证书包含两条：
\begin{enumerate}
\item 任意集合 $A$ 上有 $\mathrm{dist}(y,A)=0\iff y\in\overline{A}$；
\item 若 $A$ 闭，则 $\overline{A}=A$。
\end{enumerate}
\end{Certificate}

\begin{proof}[证明（谓词因果链）]
由证书~\ref{cert:ob-zero-criterion}：
\[
\Ob_n(y)=0
\Longleftrightarrow
y\in\overline{\mathcal A_n}
\Longleftrightarrow
y\in\mathcal A_n.
\]
命题成立。
\end{proof}

\subsection{同伦扭曲（算子参数路径）}

\begin{definition}[更新算子与同伦扭曲]\label{def:twist}
令 $R_n:\mathcal S_{n-1}\times\Theta_n\to\mathcal S_n$ 为第 $n$ 层更新（修复）算子。
给定参数序列/路径 $(\theta_n)$，定义状态序列
\[
y_n:=R_n(y_{n-1};\theta_n).
\]
把由 $\theta_n$ 的变化诱导的 $y_n$ 变形称为“同伦扭曲”。
\end{definition}

\subsection{势函数（变分泛函）}

\begin{definition}[势函数]\label{def:potential}
对每个 $n\ge3$，给定势函数（变分泛函）
\[
\Phi_n:\mathcal S_n\to\mathbb R.
\]
在构造截面或更新规则时，可采用“最小化 $\Phi_n$ 且使 $\Ob_n$ 下降”的选择策略。
\end{definition}

\section{条件化收敛：障碍修复与稳态}

\subsection{条件化收敛条款}

\begin{assumption}[障碍收缩]\label{ass:ob-contract}
存在常数序列 $(\rho_n)_{n\ge3}$ 满足 $0\le\rho_n<1$，使得沿更新 $y_n=R_n(y_{n-1};\theta_n)$ 有
\[
\Ob_n(y_n)\le \rho_n\,\Ob_{n-1}(y_{n-1}).
\]
\end{assumption}

\begin{assumption}[势值单调与下界]\label{ass:phi-monotone}
沿同一更新序列有
\[
\Phi_n(y_n)\le \Phi_{n-1}(y_{n-1}),
\]
且存在下界 $m\in\mathbb R$ 使得 $\Phi_n(y_n)\ge m$ 对所有 $n$ 成立。
\end{assumption}

\begin{theorem}[公理（降级为定理）：乘性收缩递推律（数学归纳法证书）]
\label{thm:axiom-downgrade-induction}
在假设~\ref{ass:ob-contract} 下，对任意 $n\ge3$ 有
\[
\Ob_n(y_n)
\le
\Bigl(\prod_{k=3}^{n}\rho_k\Bigr)\Ob_2(y_2).
\]
\end{theorem}

\begin{Certificate}[数学归纳法证书：按层级 $n$ 的乘积上界谓词]\label{cert:axiom-downgrade-induction}
定义谓词
\[
P(n):\ \Ob_n(y_n)\le \Bigl(\prod_{k=3}^{n}\rho_k\Bigr)\Ob_2(y_2).
\]
证书目标：
\[
P(3)\ \text{成立},\qquad
(\forall n\ge3)\ \bigl(P(n)\Rightarrow P(n+1)\bigr).
\]
\end{Certificate}

\begin{proof}[证明（数学归纳法证书；谓词因果链）]
\textbf{基步：}当 $n=3$，由假设~\ref{ass:ob-contract}
\[
\Ob_3(y_3)\le\rho_3\Ob_2(y_2)=\Bigl(\prod_{k=3}^{3}\rho_k\Bigr)\Ob_2(y_2),
\]
即 $P(3)$ 成立。

\textbf{归纳步：}设 $P(n)$ 成立。由假设~\ref{ass:ob-contract}
\[
\Ob_{n+1}(y_{n+1})\le\rho_{n+1}\Ob_n(y_n)
\le
\rho_{n+1}\Bigl(\prod_{k=3}^{n}\rho_k\Bigr)\Ob_2(y_2)
=\Bigl(\prod_{k=3}^{n+1}\rho_k\Bigr)\Ob_2(y_2).
\]
故 $P(n+1)$ 成立。

由数学归纳法，$P(n)$ 对全部 $n\ge3$ 成立。定理成立。
\end{proof}

\subsection{收敛结论}

\begin{theorem}[障碍趋零与势值收敛]\label{thm:conditional-conv}
在假设~\ref{ass:ob-contract}--\ref{ass:phi-monotone} 下，沿更新序列 $(y_n)$ 有：
\begin{enumerate}
\item 若 $\sup_{n\ge3}\rho_n\le\bar\rho<1$，则 $\Ob_n(y_n)\to0$；
\item $\Phi_n(y_n)$ 单调不增且下有界，因此存在极限
\(
\lim_{n\to\infty}\Phi_n(y_n)\in\mathbb R
\)。
\end{enumerate}
\end{theorem}

\begin{Certificate}[证书：递推上界 + 单调有界收敛]\label{cert:conditional-conv}
证书分两段：
\begin{enumerate}
\item 由定理~\ref{thm:axiom-downgrade-induction} 得乘积上界，再由
\(
\rho_n\le\bar\rho<1
\)
推出指数衰减；
\item 由假设~\ref{ass:phi-monotone} 应用实数完备性的“单调有界收敛定理”（参见 \cite{Rudin1976}）。
\end{enumerate}
\end{Certificate}

\begin{proof}[证明（谓词因果链）]
由证书~\ref{cert:conditional-conv}：
\[
\Ob_n(y_n)
\le
\Bigl(\prod_{k=3}^{n}\rho_k\Bigr)\Ob_2(y_2)
\le
\bar\rho^{n-2}\Ob_2(y_2)
\xrightarrow[n\to\infty]{}0,
\]
得第 (1) 条。
又由
\(
\Phi_n(y_n)
\)
单调不增且下有界，故极限存在，得第 (2) 条。定理成立。
\end{proof}

\begin{corollary}[任意精度可达的稳态阈值]\label{cor:epsilon-stability}
在定理~\ref{thm:conditional-conv} 的第 (1) 条件下，
对任意 $\varepsilon>0$，存在 $N(\varepsilon)$ 使得
\[
\forall n\ge N(\varepsilon),\quad \Ob_n(y_n)<\varepsilon.
\]
\end{corollary}

\begin{Certificate}[证书：收敛定义展开]\label{cert:epsilon-stability}
证书只需调用极限定义：
\(
\Ob_n(y_n)\to 0
\)
当且仅当上述 $\varepsilon$--$N$ 性质成立。
\end{Certificate}

\begin{proof}[证明（谓词因果链）]
由证书~\ref{cert:epsilon-stability} 与定理~\ref{thm:conditional-conv} 第 (1) 条直接得到结论。
\end{proof}

\begin{remark}
定理~\ref{thm:conditional-conv} 仅给出“可行性残差趋零 + 势值收敛”的稳态骨架；若要推出“收敛到全局最优”，必须额外引入下一节强条件。
\end{remark}

\section{从趋近无孔洞到存在性无孔洞：极限提升的最小接口}
\label{sec:approx-to-existence-no-hole}

\subsection{定义（把“趋近无孔洞”升级到“存在性无孔洞”所需的最小对象）}

\begin{definition}[层级塔与极限提升]\label{def:tower-limit-lift}
对每个 $n\ge 3$，给定状态空间 $\mathcal S_n$。
对每个 $n\ge 4$，给定遗忘投影
\[
\pi_{n\to n-1}:\mathcal S_n\to \mathcal S_{n-1},
\qquad
\pi_{n\to2}:=\pi_{3\to2}\circ\pi_{4\to3}\circ\cdots\circ\pi_{n\to n-1}.
\]
其中 $\pi_{3\to2}:\mathcal S_3\to\mathcal S_2$ 为底层遗忘投影。
并约定当 $n=3$ 时 $\pi_{n\to2}:=\pi_{3\to2}$，从而上述 $\pi_{n\to2}$ 可视为对全部 $n\ge3$ 的统一记号。
定义逆极限（塔极限）
\[
\mathcal S_\infty:=\varprojlim_{n\ge3}(\mathcal S_n,\pi_{n\to n-1}),
\]
其元素写作 $y_\infty=(y_3,y_4,\dots)$ 且满足一致性
\[
\pi_{n\to n-1}(y_n)=y_{n-1}\quad (n\ge4).
\]
定义极限投影 $\pi_{\infty\to2}:\mathcal S_\infty\to \mathcal S_2$ 为
\[
\pi_{\infty\to2}(y_\infty)=\pi_{3\to2}(y_3).
\]
\end{definition}

\begin{definition}[可接受子塔与极限可接受集]\label{def:acceptable-subtower-limit}
每层给定可接受集 $\mathcal A_n\subseteq \mathcal S_n$，并假设子塔相容（$n\ge4$）：
\[
\pi_{n\to n-1}(\mathcal A_n)\subseteq \mathcal A_{n-1}.
\]
定义极限可接受集
\[
\mathcal A_\infty:=\varprojlim_{n\ge3}(\mathcal A_n,\pi_{n\to n-1})\subseteq \mathcal S_\infty.
\]
对 $x\in\mathcal S_2$，定义极限提升集与极限孔洞谓词：
\[
\Lift_\infty(x):=\mathcal A_\infty\cap \pi_{\infty\to2}^{-1}(x),
\qquad
\Hole_\infty(x)\;:\Longleftrightarrow\;\Lift_\infty(x)=\varnothing.
\]
\end{definition}

\begin{definition}[上同调障碍塔与总障碍]\label{def:obstruction-total}
对每个 $k\ge2$，定义障碍类
\[
\operatorname{ob}_{k+1}(x)\in H^{k+1}(\mathcal B_x;\mathcal G_{k,x}),
\]
并定义总障碍
\[
\ObTot(x):=\bigl(\operatorname{ob}_3(x),\operatorname{ob}_4(x),\dots\bigr)\in\prod_{k\ge2} H^{k+1}(\mathcal B_x;\mathcal
  G_{k,x}).
\]
记
\[
\ObTot(x)=0\;:\Longleftrightarrow\;\forall k\ge2,\ \operatorname{ob}_{k+1}(x)=0.
\]
\end{definition}

\subsection{假设（最小完备性：把“逼近可行”升级为“存在可行”）}

\begin{assumption}[最小完备性：近似提升 $\Leftrightarrow$ 障碍归零 $\Leftrightarrow$ 极限提升存在]\label{ass:min-completeness-approx-ob-exist}
对任意 $x\in\mathcal S_2$，以下三者等价：
\begin{enumerate}
\item \textbf{存在性无孔洞（严格）：}
\(
\Lift_\infty(x)\neq\varnothing
\)；
\item \textbf{上同调障碍全归零（严格）：}
\(
\ObTot(x)=0
\)；
\item \textbf{近似无孔洞（度量意义）可实现：}存在一致的塔序列 $y_\infty=(y_3,y_4,\dots)\in\mathcal S_\infty$ 且 $\pi_{\infty\to2}(y_\infty)=x$，
  使得
\[
\Ob_n(y_n)\xrightarrow[n\to\infty]{}0,
\qquad
\Ob_n(y):=\inf_{z\in\mathcal A_n} d_n(y,z).
\]
\end{enumerate}
\end{assumption}

\begin{remark}[契约型完备性接口]
假设~\ref{ass:min-completeness-approx-ob-exist} 是一条“契约型完备性假设”：它把所有需要的拓扑/紧性/闭性/障碍理论充分性压缩为一个最小接口条款。
在不改变“两阶段结构”（先可行性修复，再最优性择优）的前提下，它把结论强度从“逼近意义的 $\Ob\to0$”升级为“存在性意义的 $\Lift_\infty(x)\neq\varnothing$”。
\end{remark}

\subsection{定理（把 $\Ob\to0$ 提升为 $\Lift_\infty(x)\neq\varnothing$）}

\begin{theorem}[严格无孔洞提升定理]\label{thm:strict-no-hole-lift-from-approx}
设第一阶段已构造出层级序列
\[
y_n = R_n(y_{n-1};\theta_n)\in\mathcal S_n\quad (n\ge3),
\]
并满足一致性与外延对齐
\[
\pi_{n\to n-1}(y_n)=y_{n-1}\quad (n\ge4),
\qquad
\pi_{n\to2}(y_n)=x\quad (n\ge3),
\]
以及逼近可行性
\[
\Ob_n(y_n)\to 0.
\]
则在假设~\ref{ass:min-completeness-approx-ob-exist} 下，
\[
\Lift_\infty(x)\neq\varnothing
\quad\text{且等价地}\quad
\ObTot(x)=0,
\]
从而 $\neg \Hole_\infty(x)$。
\end{theorem}

\begin{Certificate}[证书：假设~\ref{ass:min-completeness-approx-ob-exist} 的等价链调用]\label{cert:strict-no-hole-lift-from-approx}
证书调用假设~\ref{ass:min-completeness-approx-ob-exist} 的蕴含链
\(
\text{(3)}\Rightarrow\text{(1)}
\)
与
\(
\text{(3)}\Rightarrow\text{(2)}
\)。
\end{Certificate}

\begin{proof}[证明（谓词因果链）]
令 $y_\infty:=(y_3,y_4,\dots)$。由一致性 $\pi_{n\to n-1}(y_n)=y_{n-1}$（$n\ge4$），得 $y_\infty\in\mathcal S_\infty$，且由外延对齐得
  $\pi_{\infty\to2}(y_\infty)=x$。
又由逼近可行性 $\Ob_n(y_n)\to 0$，满足假设~\ref{ass:min-completeness-approx-ob-exist} 的条款 (3)。
因此由证书~\ref{cert:strict-no-hole-lift-from-approx} 得
\(
\Lift_\infty(x)\neq\varnothing
\)
与
\(
\ObTot(x)=0
\)。
再由定义~\ref{def:acceptable-subtower-limit} 中
\(
\Hole_\infty(x)\Longleftrightarrow \Lift_\infty(x)=\varnothing
\)
知 $\neg \Hole_\infty(x)$。定理成立。
\end{proof}

\subsection{推论（把“孔洞”与“上同调障碍归零”严格等价化）}

\begin{corollary}[严格充要式：孔洞 $\Leftrightarrow$ 障碍非零]\label{cor:hole-iff-obtot-strict}
在假设~\ref{ass:min-completeness-approx-ob-exist} 下，对任意 $x\in\mathcal S_2$，有严格充要式：
\[
\Hole_\infty(x)
\Longleftrightarrow
\Lift_\infty(x)=\varnothing
\Longleftrightarrow
\ObTot(x)\neq 0.
\]
\end{corollary}

\begin{Certificate}[证书：由“$\Lift_\infty\neq\varnothing \Longleftrightarrow \ObTot=0$”取否定]\label{cert:hole-iff-obtot-strict}
证书包含两步：
\begin{enumerate}
\item 由定义~\ref{def:acceptable-subtower-limit} 得 $\Hole_\infty(x)\Longleftrightarrow \Lift_\infty(x)=\varnothing$；
\item 由假设~\ref{ass:min-completeness-approx-ob-exist} 的 (1)$\Longleftrightarrow$(2) 得 $\Lift_\infty(x)
  \neq\varnothing\Longleftrightarrow \ObTot(x)=0$，两侧同取否定得到 $\Lift_\infty(x)=\varnothing\Longleftrightarrow \ObTot(x)\neq
  0$。
\end{enumerate}
\end{Certificate}

\begin{proof}[证明（谓词因果链）]
由证书~\ref{cert:hole-iff-obtot-strict} 两步串接，立即得到所述三段等价。推论成立。
\end{proof}

\subsection{备注（关于“截面”的一句话，避免误解）}

\begin{remark}[存在性无孔洞与极限截面]
\begin{enumerate}
\item 本次升级的核心是把“$\Ob\to0$”从\textbf{逼近意义}升级为“$\Lift_\infty(x)\neq\varnothing$”的\textbf{存在性意义}。这一步不需要改变主定理的两段结构（先可行性修复，
  再最优性择优），只需补上假设~\ref{ass:min-completeness-approx-ob-exist}。
\item 若希望从“对每个 $x$ 有提升”进一步升级到“存在一个截面 $\sigma_\infty$”，例如
\[
\sigma_\infty:\mathcal D_\infty\to\mathcal A_\infty,\quad
\pi_{\infty\to2}\circ\sigma_\infty=\mathrm{id}_{\mathcal D_\infty},
\quad
\mathcal D_\infty:=\{x\in\mathcal S_2:\Lift_\infty(x)\neq\varnothing\},
\]
则通常只需集合论层面的“选择”（对每个 $x$ 选取一个提升代表）。该截面是否\textbf{可构造}属于额外要求；不加则仍是严格的“存在性无孔洞”。
\end{enumerate}
\end{remark}

\section{变分法：全局最优存在性与全局收敛条件}

\subsection{最优解存在性（可达性）}

\begin{theorem}[最小值可达（存在性条款）]\label{thm:existence-minimizer}
对固定外延态 $x\in\mathcal S_2$，令可行集
\[
K(x):=\Lift_\infty(x).
\]
若满足：
\begin{enumerate}
\item $K(x)\neq\varnothing$；
\item $K(x)$ 在所用拓扑下为闭集；
\item $\inf_{y\in K(x)}\Phi(y)>-\infty$；
\item $K(x)$ 紧，或 $\Phi$ 在 $K(x)$ 上胁迫（当 $\|y\|\to\infty$ 且 $y\in K(x)$ 时 $\Phi(y)\to+\infty$），
\end{enumerate}
则存在
\(
y^\star(x)\in\arg\min_{y\in K(x)}\Phi(y)
\)。
\end{theorem}

\begin{Certificate}[证书：极小化序列紧化与闭性回收]\label{cert:existence-minimizer}
证书链条：
\[
\begin{aligned}
\text{下有界}
&\Rightarrow \text{存在极小化序列}
\Rightarrow \text{紧性/胁迫性给出有界并可抽收敛子列}\\
&\Rightarrow \text{闭性保证极限留在 }K(x)
\Rightarrow \text{极限达到下确界}.
\end{aligned}
\]
\end{Certificate}

\begin{remark}[参考：直接法/存在性]
定理~\ref{thm:existence-minimizer} 的证明骨架可视为变分法“直接法”在抽象可行集上的模板化改写，可参见 \cite{Dacorogna2008}。
\end{remark}

\begin{proof}[证明（谓词因果链）]
由证书~\ref{cert:existence-minimizer}：取极小化序列 $(y^{(m)})\subseteq K(x)$，并令其满足
\[
\Phi(y^{(m)})\downarrow\inf_{y\in K(x)}\Phi(y).
\]
由紧性或胁迫性，可抽取收敛子列 $y^{(m_j)}$，使 $y^{(m_j)}\to y^\star$。
闭性给出 $y^\star\in K(x)$。
再由下半连续型极限传递（在本文口径下由构造假设吸收）得
\[
\Phi(y^\star)=\inf_{y\in K(x)}\Phi(y).
\]
故最小值可达。
\end{proof}

\subsection{收敛到全局最优（强条件）}

\begin{assumption}[凸性与强凸性]\label{ass:strong-convex}
假设 $K(x)$ 为凸集，且 $\Phi$ 在 $K(x)$ 上 $\mu$-强凸：存在 $\mu>0$ 使
\[
\Phi(v)\ge \Phi(u)+\langle\nabla\Phi(u),v-u\rangle+\frac{\mu}{2}\|v-u\|^2,
\quad \forall u,v\in K(x).
\]
\end{assumption}

\begin{assumption}[Lipschitz 梯度与步长]\label{ass:lipschitz-step}
假设 $\nabla\Phi$ 在 $K(x)$ 上 $L$-Lipschitz，且取步长 $0<\alpha<2/L$。
\end{assumption}

\begin{theorem}[投影梯度法全局收敛（强条件版）]\label{thm:pgd-global}
在假设~\ref{ass:strong-convex}--\ref{ass:lipschitz-step} 下，迭代
\[
y_{t+1}=\Pi_{K(x)}\bigl(y_t-\alpha\nabla\Phi(y_t)\bigr)
\]
收敛到唯一全局最优解 $y^\star(x)$。
\end{theorem}

\begin{Certificate}[证书：强凸唯一性 + 非扩张投影 + 梯度步收缩]\label{cert:pgd-global}
证书登记三项标准接口：
\begin{enumerate}
\item 强凸性 $\Rightarrow$ 最优解唯一；
\item 投影算子 $\Pi_{K(x)}$ 为非扩张映射；
\item 在 $0<\alpha<2/L$ 下，梯度步与强凸组合给出误差收缩。
\end{enumerate}
\end{Certificate}

\begin{remark}[参考：凸优化与一阶方法]
投影梯度法在强凸 + Lipschitz 梯度条件下的收敛性属于凸优化与一阶方法的标准结论，可参见 \cite{BoydVandenberghe2004,Nesterov2004,Beck2017,BauschkeCombettes2017}
  。
\end{remark}

\begin{proof}[证明（谓词因果链）]
由证书~\ref{cert:pgd-global}：
\[
\text{强凸}
\Rightarrow
\exists!\ y^\star(x),
\]
且对误差 $e_t:=y_t-y^\star$ 有
\[
\|e_{t+1}\|
\le
q\|e_t\|,\qquad 0<q<1.
\]
于是
\(
\|e_t\|\le q^t\|e_0\|\to0
\)
，故 $y_t\to y^\star(x)$。定理成立。
\end{proof}

\begin{corollary}[强条件下的唯一极小值判据]\label{cor:unique-minimizer}
在假设~\ref{ass:strong-convex} 下，$K(x)$ 上最小值点唯一。
\end{corollary}

\begin{Certificate}[证书：强凸排斥双极小点]\label{cert:unique-minimizer}
若存在 $u\neq v$ 均为极小点，代入强凸不等式并取中点
\(
\tfrac{u+v}{2}
\)
会得到严格更小函数值，矛盾。
\end{Certificate}

\begin{proof}[证明（谓词因果链）]
由证书~\ref{cert:unique-minimizer}：假设存在两极小点 $u\neq v$，则强凸性推出
\[
\Phi\Bigl(\frac{u+v}{2}\Bigr)
<
\frac{\Phi(u)+\Phi(v)}{2}
=
\Phi(u),
\]
与“$u$ 为极小点”矛盾。故极小点唯一。
\end{proof}

\begin{remark}
若缺乏凸性/强凸性/Lipschitz 梯度等条件，则“全局最优路径收敛”一般属于\textbf{未证明}；此时应回退到第~\ref{thm:conditional-conv} 的稳态结论。
\end{remark}

\section{双重收敛：层级修复与可行域优化的可审计分解}
\label{sec:double-convergence}

\subsection{记号与目标}
\label{subsec:double-convergence-notation-goal}

以下严格区分两类索引：
\begin{itemize}
\item \textbf{层级索引：}$n\ge 3$，表示语义层 $l_n$ 的层级修复（同伦/曲率/语义增强的层）；对应层级更新序列 $y_n$；
\item \textbf{优化迭代索引：}$t\in\mathbb N$，表示在固定可行域 $K(x)$ 上做变分迭代（投影梯度步）；对应优化迭代序列 $y^{t}$。
\end{itemize}

本文把“收敛”严格拆成两条可审计结论，并强调二者的接口分工：
\begin{enumerate}
\item \textbf{可行性收敛（层级修复轨道）}：沿层级更新 $y_n$，障碍量趋零 $\Ob_n(y_n)\to 0$；
\item \textbf{最优性收敛（可行域内择优轨道）}：在可行域 $K(x)$ 上，沿优化迭代 $y^{t}$ 收敛到唯一全局最优 $y^\star(x)$。
\end{enumerate}

\begin{remark}[接口化直觉]
第一条只管“是否逼近可接受集”（一致性/可行性）；第二条只管“在可行域内如何择优”（全局最优/收敛率）。
两者相互正交：$\Ob$ 管可行性残差，$\Phi$ 管最优性比较。
\end{remark}

\subsection{定义（层级数据、障碍量、提升与最优）}
\label{subsec:double-convergence-definitions}

\begin{definition}[D1（层级数据与语义增强）]\label{def:dc-D1}
对每个 $n\ge 3$，定义层级数据五元组
\[
\mathfrak D_n := \bigl(\mathcal S_n,\ \Theta_n,\ R_n,\ d_n,\ \mathcal A_n\bigr),
\]
其中：$\mathcal S_n$ 为第 $n$ 层状态空间（见定义~\ref{def:layer-space}），$\Theta_n$ 为连续参数空间（见定义~\ref{def:mixed-index}），
$R_n:\mathcal S_{n-1}\times \Theta_n\to \mathcal S_n$ 为更新/修复算子（见定义~\ref{def:twist}），
$d_n$ 为 $\mathcal S_n$ 上度量（或伪度量，见定义~\ref{def:ob-measure}），
$\mathcal A_n\subseteq \mathcal S_n$ 为可接受集（见定义~\ref{def:acceptable}）。
\end{definition}

\begin{remark}[对比：$l_2$ 的接口缺失]
层 $l_2$ 仅保留外延信息（离散标签/基数级）。在不额外引入 $d_n,\mathcal A_n,\Phi$ 等接口的前提下，
无法在 $l_2$ 内部形成可审计的“修复收敛”与“最优收敛”不等式链。
\end{remark}

\begin{definition}[D2（障碍量）]\label{def:dc-D2}
定义障碍量（到可接受集的距离）
\[
\Ob_n(y):=\inf_{z\in\mathcal A_n} d_n(y,z)\in\mathbb R_{\ge 0}
\]
（见定义~\ref{def:ob-measure}）。它刻画“可行性残差”（一致性缺口/孔洞风险的度量化替身）。
\end{definition}

\begin{definition}[D3（纤维、提升、孔洞）]\label{def:dc-D3}
给定遗忘投影
\[
\pi_{n\to 2}:\mathcal S_n\to\mathcal S_2
\]
（见定义~\ref{def:forget-map}）。对任意 $x\in\mathcal S_2$ 定义可接受提升集
\[
\Lift_n(x):=\mathcal A_n\cap \pi_{n\to 2}^{-1}(x)
\]
（见定义~\ref{def:lift}），并定义孔洞谓词
\[
\Hole_n(x)\ :\iff\ \Lift_n(x)=\varnothing
\]
（见定义~\ref{def:hole}）。
\end{definition}

\begin{remark}[从 $\Ob\to 0$ 到“无孔洞”需要额外条款]
本章“可行性收敛”的主结论以 $\Ob_n(y_n)\to 0$ 表达。若要从 $\Ob_n(y_n)\to 0$ 进一步推出
“确实消除孔洞”（例如 $x$ 具有极限意义下的可接受提升），需要额外加入“提升完备性/极限存在性/障碍完备性”等条款，
参见第~\ref{subsec:double-convergence-remarks} 小节。
\end{remark}

\begin{definition}[D4（可行域与势函数）]\label{def:dc-D4}
对固定外延态 $x\in\mathcal S_2$，令 $K(x)$ 为承载优化的目标可行域（可取 $K(x)=\Lift_\infty(x)$，见定理~\ref{thm:existence-minimizer} 的口径）。
给定势函数（变分泛函）
\[
\Phi:K(x)\to\mathbb R,
\]
并定义全局最优解
\[
y^\star(x)\in \arg\min_{y\in K(x)}\Phi(y).
\]
\end{definition}

\subsection{公理/假设（证书）}
\label{subsec:double-convergence-assumptions}

\begin{assumption}[障碍收缩（可行性修复）]\label{ass:dc-ob-contract}
沿层级更新序列
\[
y_n = R_n(y_{n-1};\theta_n),\qquad n\ge 3,
\]
存在 $\rho_n\in[0,1)$ 使得
\[
\Ob_n(y_n)\le \rho_n\,\Ob_{n-1}(y_{n-1}),
\]
并且存在统一上界 $\bar\rho<1$ 满足
\[
\sup_{n\ge 3}\rho_n \le \bar\rho.
\]
\end{assumption}

\begin{assumption}[强凸 + Lipschitz 梯度（最优性锁定）]\label{ass:dc-opt-lock}
在可行域 $K(x)$ 上满足：
\begin{enumerate}
\item $K(x)$ 非空、闭、凸；
\item $\Phi$ 为 $\mu$-强凸（$\mu>0$），即
\[
\Phi(v)\ge \Phi(u)+\langle \nabla\Phi(u),v-u\rangle + \frac{\mu}{2}\|v-u\|^2,\quad \forall u,v\in K(x);
\]
\item $\nabla \Phi$ 为 $L$-Lipschitz（$L>0$），即
\[
\|\nabla\Phi(u)-\nabla\Phi(v)\|\le L\|u-v\|,\quad \forall u,v\in K(x).
\]
\end{enumerate}
\end{assumption}

\begin{Certificate}[证书 0：接口引用（结构前提）]\label{cert:dc-interface}
本章将以下接口作为可直接调用的结构前提（无需在本文内部重证）：
\begin{itemize}
\item 修复算子接口：$R_n:\mathcal S_{n-1}\times\Theta_n\to\mathcal S_n$；
\item 障碍量定义：$\Ob_n(y)=\inf_{z\in\mathcal A_n}d_n(y,z)$；
\item 投影算子：$\Pi_{K(x)}$ 为到闭凸集 $K(x)$ 的投影（因此满足非扩张性；见引理~\ref{lem:dc-proj-nonexpansive}，参考 \cite{BauschkeCombettes2017}）；
\item 优化迭代（投影梯度）：
\[
y^{t+1}=\Pi_{K(x)}\bigl(y^t-\alpha\nabla\Phi(y^t)\bigr).
\]
\end{itemize}
\end{Certificate}

\subsection{第一重收敛：可行性修复（障碍量归零）}
\label{subsec:double-convergence-feasibility}

\begin{lemma}[障碍量指数衰减]\label{lem:dc-ob-decay}
在假设~\ref{ass:dc-ob-contract} 下，
\[
\Ob_n(y_n)
\le
\Bigl(\prod_{k=3}^n \rho_k\Bigr)\Ob_2(y_2)
\le
\bar\rho^{\,n-2}\Ob_2(y_2)
\xrightarrow[n\to\infty]{}0.
\]
\end{lemma}

\begin{Certificate}[数学归纳法证书：乘积上界谓词]\label{cert:dc-ob-decay}
令命题
\[
P(n):\quad \Ob_n(y_n)\le\Bigl(\prod_{k=3}^n\rho_k\Bigr)\Ob_2(y_2).
\]
证书目标：
\[
P(3)\ \text{成立},\qquad (\forall n\ge 3)\ \bigl(P(n)\Rightarrow P(n+1)\bigr).
\]
并在归纳结论基础上调用 $\rho_k\le\bar\rho<1$ 推出极限为 $0$。
\end{Certificate}

\begin{proof}[证明（数学归纳法证书；谓词因果链）]
\textbf{基步：}当 $n=3$，由假设~\ref{ass:dc-ob-contract} 得
\[
\Ob_3(y_3)\le\rho_3\Ob_2(y_2)=\Bigl(\prod_{k=3}^{3}\rho_k\Bigr)\Ob_2(y_2),
\]
故 $P(3)$ 成立。

\textbf{归纳步：}设 $P(n)$ 成立。由假设~\ref{ass:dc-ob-contract} 得
\[
\Ob_{n+1}(y_{n+1})\le \rho_{n+1}\Ob_n(y_n)
\le \rho_{n+1}\Bigl(\prod_{k=3}^n\rho_k\Bigr)\Ob_2(y_2)
=\Bigl(\prod_{k=3}^{n+1}\rho_k\Bigr)\Ob_2(y_2),
\]
故 $P(n+1)$ 成立。

由数学归纳法得 $P(n)$ 对全部 $n\ge 3$ 成立。再由 $\rho_k\le\bar\rho<1$ 得
\[
\Ob_n(y_n)\le\bar\rho^{\,n-2}\Ob_2(y_2)\xrightarrow[n\to\infty]{}0.
\]
引理成立。
\end{proof}

\begin{corollary}[趋于可接受：任意精度可行]\label{cor:dc-eps-feasible}
在引理~\ref{lem:dc-ob-decay} 下，对任意 $\varepsilon>0$，存在 $N(\varepsilon)$ 使得
\[
\forall n\ge N(\varepsilon),\quad \Ob_n(y_n)\le \varepsilon.
\]
\end{corollary}

\begin{Certificate}[证书：极限定义展开]\label{cert:dc-eps-feasible}
证书只需把 $\Ob_n(y_n)\to 0$ 按 $\varepsilon$--$N$ 语言展开即可。
\end{Certificate}

\begin{proof}[证明（谓词因果链）]
由证书~\ref{cert:dc-eps-feasible} 与引理~\ref{lem:dc-ob-decay} 直接推出结论。
\end{proof}

\subsection{第二重收敛：最优路径锁定（全局最优）}
\label{subsec:double-convergence-optimality}

\begin{lemma}[投影非扩张性]\label{lem:dc-proj-nonexpansive}
对闭凸集 $K(x)$ 的投影 $\Pi_{K(x)}$，有
\[
\|\Pi_{K(x)}(u)-\Pi_{K(x)}(v)\|\le \|u-v\|,\quad \forall u,v.
\]
\end{lemma}

\begin{Certificate}[证书：投影的一阶最优性条件]\label{cert:dc-proj-nonexpansive}
证书登记：闭凸集上的投影满足变分不等式刻画，并推出（甚至更强的）firmly nonexpansive 性质；可参见 \cite{BauschkeCombettes2017}。
\end{Certificate}

\begin{proof}[证明（谓词因果链）]
由证书~\ref{cert:dc-proj-nonexpansive}：
\[
\text{闭凸集}\Rightarrow \Pi_{K(x)}\ \text{存在且唯一}
\Rightarrow \Pi_{K(x)}\ \text{满足投影变分不等式}
\Rightarrow \Pi_{K(x)}\ \text{非扩张}.
\]
引理成立。
\end{proof}

\begin{proposition}[唯一全局最优解存在]\label{prop:dc-unique-opt}
在假设~\ref{ass:dc-opt-lock} 下，$\Phi$ 在 $K(x)$ 上存在唯一全局最优解 $y^\star(x)$。
\end{proposition}

\begin{Certificate}[证书：强凸唯一性 + 胁迫性/可达性模板]\label{cert:dc-unique-opt}
证书链条：
\[
\text{强凸}\Rightarrow \text{极小点至多一个}
\Rightarrow \text{唯一性成立};
\]
并在本文默认口径下，将 $K(x)$ 视为欧氏空间上的闭凸集，则强凸性给出二次下界，从而 $\Phi$ 在 $K(x)$ 上胁迫，
结合闭性与连续性得到极小值可达（参见 \cite{BoydVandenberghe2004,Beck2017}）。
\end{Certificate}

\begin{proof}[证明（谓词因果链）]
由证书~\ref{cert:dc-unique-opt}：
\[
\text{强凸}\Rightarrow \text{若存在两极小点则中点严格更小}\Rightarrow \text{矛盾}\Rightarrow \text{极小点唯一}.
\]
另一方面，强凸性给出二次增长下界，从而 $\Phi$ 在 $K(x)$ 上胁迫；结合 $K(x)$ 的闭性与 $\Phi$ 的连续性，下确界可达。
故唯一最小点存在。命题成立。
\end{proof}

\begin{proposition}[投影梯度线性收敛]\label{prop:dc-pgd-linear}
在假设~\ref{ass:dc-opt-lock} 下，取步长满足
\[
0<\alpha<\frac{2}{L},
\]
并定义迭代
\[
y^{t+1}=\Pi_{K(x)}\bigl(y^t-\alpha\nabla\Phi(y^t)\bigr).
\]
则存在 $q\in(0,1)$ 使得
\[
\|y^{t+1}-y^\star(x)\|\le q\,\|y^{t}-y^\star(x)\|,\qquad \forall t,
\]
从而 $y^{t}\to y^\star(x)$（线性收敛）。
\end{proposition}

\begin{Certificate}[证书：压缩映射 + 投影非扩张 + 不动点条件]\label{cert:dc-pgd-linear}
证书分三步：
\begin{enumerate}
\item 对 $\mu$-强凸且 $L$-smooth 的 $\Phi$，在合适步长下，梯度步进映射 $T(y):=y-\alpha\nabla\Phi(y)$ 在 $K(x)$ 上为压缩映射；
\item 由引理~\ref{lem:dc-proj-nonexpansive}，投影 $\Pi_{K(x)}$ 为非扩张映射；
\item 最优解 $y^\star$ 满足投影梯度不动点条件 $y^\star=\Pi_{K(x)}(T(y^\star))$（等价于闭凸集上一阶最优性）。
\end{enumerate}
上述为一阶方法标准链，可参见 \cite{Nesterov2004,Beck2017,BauschkeCombettes2017}。
\end{Certificate}

\begin{proof}[证明（谓词因果链）]
由证书~\ref{cert:dc-pgd-linear}：
\[
\|y^{t+1}-y^\star\|
=\|\Pi_{K(x)}(T(y^t))-\Pi_{K(x)}(T(y^\star))\|
\le \|T(y^t)-T(y^\star)\|
\le q\,\|y^t-y^\star\|,
\]
其中 $0<q<1$。故误差满足线性收缩，从而 $y^t\to y^\star(x)$。命题成立。
\end{proof}

\begin{remark}[边界提示]
若去掉强凸/闭凸/步长等条件，“收敛到全局最优”一般属于\textbf{未证明}；通常只能证明“势值收敛到某个临界值或停点”。
\end{remark}

\subsection{核心定理：$l_{\ge 3}$ 的双重收敛修复（可行性 + 最优性）}
\label{subsec:double-convergence-main-theorem}

\begin{theorem}[双重收敛主定理]\label{thm:dc-double-convergence}
在 $n\ge 3$ 的语义层级上，若满足：
\begin{itemize}
\item 假设~\ref{ass:dc-ob-contract}（障碍收缩：可行性修复）；
\item 存在非空闭凸可行域 $K(x)$ 与势函数 $\Phi$ 满足假设~\ref{ass:dc-opt-lock}（强凸 + Lipschitz 梯度：最优性锁定）；
\item 并在可行域 $K(x)$ 上以投影梯度法进行优化迭代；
\end{itemize}
则存在两重收敛结论：
\begin{enumerate}
\item \textbf{(A) 可行性修复收敛：}沿层级修复序列 $y_n$，有 $\Ob_n(y_n)\to 0$；
\item \textbf{(B) 最优路径锁定收敛：}沿优化迭代 $y^t$，有 $y^t\to y^\star(x)$，其中 $y^\star(x)$ 为 $K(x)$ 上 $\Phi$ 的唯一全局最优解。
\end{enumerate}
\end{theorem}

\begin{Certificate}[证书：两条收敛链的拼接]\label{cert:dc-double-convergence}
证书链条：
\[
\text{假设~\ref{ass:dc-ob-contract}}
\Rightarrow
\text{引理~\ref{lem:dc-ob-decay}}
\Rightarrow
\text{(A)};
\]
\[
\text{假设~\ref{ass:dc-opt-lock}}
\Rightarrow
\text{命题~\ref{prop:dc-unique-opt} 与命题~\ref{prop:dc-pgd-linear}}
\Rightarrow
\text{(B)}.
\]
两条链共享的仅是“接口对象的可引用性”（证书~\ref{cert:dc-interface}），在逻辑上彼此独立。
\end{Certificate}

\begin{proof}[证明（谓词因果链）]
由证书~\ref{cert:dc-double-convergence}：
(A) 由引理~\ref{lem:dc-ob-decay} 直接得到；
(B) 由命题~\ref{prop:dc-unique-opt} 得最优解存在且唯一，再由命题~\ref{prop:dc-pgd-linear} 得线性收敛到该唯一最优解。
两者合并即得结论。定理成立。
\end{proof}

\subsection{备注：孔洞/上同调障碍与双重收敛对齐}
\label{subsec:double-convergence-remarks}

\begin{remark}
为把“孔洞/上同调障碍”与本章双重收敛结论对齐，需要区分“逼近”与“存在”两类语义强度：
\begin{enumerate}
\item \textbf{从孔洞到无孔洞的最弱落地：}本章把“从孔洞到无孔洞”的最弱可审计表达落实为 $\Ob_n(y_n)\to 0$，
它等价于“以任意精度逼近可接受集”（推论~\ref{cor:dc-eps-feasible}）。
\item \textbf{从 $\Ob\to 0$ 升级到“确有提升/确有截面”：}若要求真正得到 $x$ 的极限意义下的可接受提升（或截面存在），可加入例如：
\begin{itemize}
\item 障碍完备性：$\Lift_\infty(x)\neq\varnothing \iff \ObTot(x)=0$（假设~\ref{ass:ot-complete} 与定理~\ref{thm:hole-iff-obtot}）；
\item 极限存在性：存在极限对象 $y_\infty\in K(x)$ 且满足某种极限投影条件（例如 $\pi_{\infty\to 2}(y_\infty)=x$）。
\end{itemize}
\item \textbf{$l_2$ 的局限：}缺少 $(d,\Ob,\Phi)$ 等接口时，“收缩/优化”无法落到不等式与收敛率；因此 $l_2$ 只能呈现外延标签，而不能自产生“修复--择优”的证明链。
\item \textbf{两种无穷（可数 + 不可数）：}层级索引 $n\in\mathbb N$（可数）与连续参数 $\theta_n\in\Theta_n$（不可数）共同构成修复自由度；
这正是让“绕过障碍”与“势能比较”成为可计算/可审计机制的关键结构来源（见定义~\ref{def:mixed-index} 与引理~\ref{lem:two-infinities}）。
\end{enumerate}
\end{remark}

\section{两类无穷自由度：全连通与最优路径的严格化表述}
\label{sec:two-infinities-connectivity-optimal}

\subsection{定义}

\begin{definition}[D1（两类“无穷自由度”）]\label{def:D1-two-infinities-freedom}
本文中与“修复--择优”相关的自由度分为两类：
\begin{enumerate}
\item \textbf{可数无穷自由度：层级/塔阶索引}
\[
n\in\mathbb N_{\ge3},
\]
表示“第 $n$ 级新增一批边缘算子/修复算子”，用于逐级处理障碍塔 $\bigl(\operatorname{ob}_{k+1}(x)\bigr)_{k\ge2}$（见定义~\ref{def:obstruction-tower}
  与定义~\ref{def:obstruction-total}）。
\item \textbf{不可数无穷自由度：每级内部的连续语义参数}
\[
\theta_n\in\Theta_n,\qquad |\Theta_n|\ge|\mathbb R|,
\]
用于在固定可行约束下引入连续变形、梯度与势函数比较结构，从而执行变分择优（见定义~\ref{def:mixed-index} 与定义~\ref{def:potential}）。
\end{enumerate}
\end{definition}

\subsection{命题}

\begin{proposition}[P1（可数无穷自由度“实现全连通”的严格化版本）]\label{prop:P1-countable-global-feasible-connectivity}
把“全连通”严格化为以下更可审计的目标：
\[
\textbf{全连通（可行连通）}:=\forall x\in\mathcal S_2,\ \Lift_\infty(x)\neq\varnothing
\quad(\text{即 }\Hole_\infty(x)\ \text{处处为假}).
\]
在假设~\ref{ass:min-completeness-approx-ob-exist} 下，若对任意 $x\in\mathcal S_2$ 都存在可数层级更新序列
\[
y_n=R_n(y_{n-1};\theta_n)\in\mathcal S_n\quad(n\ge3),
\qquad
\pi_{n\to2}(y_n)=x,
\]
并且满足
\[
\Ob_n(y_n)\xrightarrow[n\to\infty]{}0,
\]
则推出 $\Lift_\infty(x)\neq\varnothing$，从而“无孔洞/可行连通”成立。
\end{proposition}

\begin{Certificate}[证书：收缩趋零接口 + 最小完备性升级接口]\label{cert:P1-countable-global-feasible-connectivity}
证书分两段：
\begin{enumerate}
\item \textbf{（可选的趋零来源）}若沿层级更新满足假设~\ref{ass:ob-contract}，并进一步满足定理~\ref{thm:conditional-conv} 的条件 $\sup_{n\ge3}
  \rho_n\le\bar\rho<1$，则可推出 $\Ob_n(y_n)\to0$；
\item \textbf{（结论升级）}由假设~\ref{ass:min-completeness-approx-ob-exist} 的 (3)$\Rightarrow$(1) 与 (3)$\Rightarrow$(2)
  （或直接调用定理~\ref{thm:strict-no-hole-lift-from-approx}），把 $\Ob_n(y_n)\to0$ 升级为 $\Lift_\infty(x)\neq\varnothing$（并推出
  $\ObTot(x)=0$）。
\end{enumerate}
\end{Certificate}

\begin{proof}[证明（谓词因果链）]
任取 $x\in\mathcal S_2$。由命题条件，存在一致层级更新序列 $(y_n)_{n\ge3}$ 满足 $\pi_{n\to2}(y_n)=x$ 且 $\Ob_n(y_n)\to0$。
因此满足假设~\ref{ass:min-completeness-approx-ob-exist} 的条款 (3)。
由证书~\ref{cert:P1-countable-global-feasible-connectivity} 的第 2 段，得到 $\Lift_\infty(x)\neq\varnothing$（并且 $\ObTot(x)=0$）。
由于 $x$ 任意，故 $\forall x\in\mathcal S_2,\ \Lift_\infty(x)\neq\varnothing$，即“全连通（可行连通）”成立。命题成立。
\end{proof}

\subsection{定理}

\begin{theorem}[T1（不可数无穷自由度“实现最优路径”的严格化版本）]\label{thm:T1-uncountable-optimal-path}
对固定 $x\in\mathcal S_2$，设可行域 $K(x)\neq\varnothing$，并在本文口径下取
\[
K(x):=\Lift_\infty(x).
\]
若满足：
\begin{enumerate}
\item $K(x)$ 闭凸；
\item $\Phi$ 在 $K(x)$ 上 $\mu$-强凸，且 $\nabla\Phi$ 在 $K(x)$ 上为 $L$-Lipschitz；
\end{enumerate}
则投影梯度迭代
\[
y^{t+1}=\Pi_{K(x)}\bigl(y^t-\alpha\nabla\Phi(y^t)\bigr),\qquad 0<\alpha<\frac{2}{L}
\]
收敛到唯一全局最优解 $y^\star(x)$，并且存在 $q\in(0,1)$ 使误差线性收缩：
\[
\|y^{t+1}-y^\star(x)\|\le q\,\|y^t-y^\star(x)\|,\quad \forall t.
\]
\end{theorem}

\begin{Certificate}[证书：强凸唯一性 + 压缩映射 + 投影非扩张]\label{cert:T1-uncountable-optimal-path}
证书登记三项标准接口：
\begin{enumerate}
\item 强凸 + Lipschitz 梯度 $\Rightarrow$ 全局最优解唯一；
\item 在 $0<\alpha<2/L$ 下，梯度步进映射为压缩；
\item 闭凸集投影 $\Pi_{K(x)}$ 非扩张，与压缩步进复合仍为压缩，从而得到线性收敛。
\end{enumerate}
上述链条在本文中可直接由命题~\ref{prop:dc-unique-opt} 与命题~\ref{prop:dc-pgd-linear} 调用。
\end{Certificate}

\begin{proof}[证明（谓词因果链）]
由证书~\ref{cert:T1-uncountable-optimal-path}：在给定的闭凸可行域 $K(x)$ 上，
\[
\begin{aligned}
\text{强凸 + Lipschitz 梯度}
&\Rightarrow
\exists!\ y^\star(x)\\
&\Rightarrow
\text{投影梯度迭代为压缩}\\
&\Rightarrow
\exists q\in(0,1)\ \text{使得}\\
&\qquad
\|y^{t+1}-y^\star\|\le q\,\|y^t-y^\star\|\\
&\Rightarrow
y^t\to y^\star(x).
\end{aligned}
\]
严格推导见命题~\ref{prop:dc-unique-opt} 与命题~\ref{prop:dc-pgd-linear}。定理成立。
\end{proof}

\subsection{推论}

\begin{corollary}[C1（“可数无穷实现全连通，不可数无穷实现最优路径”的严格版本）]\label{cor:C1-two-infinities-slogan}
您所述句子
\[
\text{“可数无穷（自由度）实现全连通，不可数无穷（自由度）实现最优路径”}
\]
在本文框架下可严格改写为两条可审计断言：
\[
\boxed{
\begin{aligned}
\text{可数层级 }(n\in\mathbb N)\ \text{用于逐级消除障碍残差}
&\Rightarrow
\forall x\in\mathcal S_2,\ \Lift_\infty(x)\neq\varnothing\\
&\qquad(\text{无孔洞/可行连通}).
\end{aligned}
}
\]
\[
\boxed{
\begin{aligned}
\text{不可数连续参数 }(\theta_n\in\Theta_n,\ |\Theta_n|\ge|\mathbb R|)\ \text{支撑势函数与梯度结构}
&\Rightarrow
y^t\to y^\star(x)\\
&\qquad(\text{最优路径收敛}).
\end{aligned}
}
\]
\end{corollary}

\begin{Certificate}[证书：命题 P1 + 定理 T1 直接拼接]\label{cert:C1-two-infinities-slogan}
证书仅需两项：
\begin{enumerate}
\item 由命题~\ref{prop:P1-countable-global-feasible-connectivity} 得“无孔洞/可行连通”；
\item 由定理~\ref{thm:T1-uncountable-optimal-path} 得“最优路径收敛”。
\end{enumerate}
\end{Certificate}

\begin{proof}[证明（谓词因果链）]
由证书~\ref{cert:C1-two-infinities-slogan}，两条方框断言分别是命题~\ref{prop:P1-countable-global-feasible-connectivity} 与定理~\ref{thm:T1-uncountable-optimal-path} 的文字化重述。推论成立。
\end{proof}

\subsection{备注（两点必要的精确化）}

\begin{remark}[全连通与不可数自由度的边界]\label{rem:two-infinities-precision}
\begin{enumerate}
\item “全连通”若按拓扑学标准理解为“路径连通/单连通”等更强性质，则需额外假设（例如各层可接受集 $\mathcal A_n$ 的路径连通性，以及投影/提升保持连通性）。
在目前最小骨架里，“全连通”最稳妥的严格含义是\textbf{处处无孔洞（可行提升处处存在）}，即 $\forall x,\ \Lift_\infty(x)\neq\varnothing$。
\item “不可数”并非逻辑上绝对必要：离散优化同样可以寻优。
但在本文口径下，“不可数自由度”的关键作用是为“势函数比较、梯度、变分路径、收缩率”等\textbf{可审计不等式}提供自然载体，从而把“最优路径”从叙述目标提升为可证明对象。
\end{enumerate}
\end{remark}

\section{两条链的严格对齐：Chain-B 作为 Chain-A 的实例}

\subsection{对齐主定理}

\begin{theorem}[链对齐]\label{thm:alignment}
在前述对象与假设下，Chain-B 可按如下对应实现为 Chain-A 的实例化：
\[
\begin{array}{rcl}
l_2 &\leftrightarrow& \mathcal S_2,\\
\text{连续统（一离散）} &\leftrightarrow& \text{纤维 }\operatorname{Fib}_n(x)\text{ 的冗余 + }l_2\text{ 的离散外延输出},\\
\text{崩溃/上同调障碍} &\leftrightarrow& \Hole_n(x)\ (\text{并在 }\mathbf{(OT\text{-}Complete)}\text{ 下等价于 }\ObTot(x)\neq0),\\
\text{同伦塔度量（}l_{\ge3}\text{）} &\leftrightarrow& (\mathcal S_n,d_n)\text{ 与 }\Ob_n,\\
\text{同伦扭曲} &\leftrightarrow& y_n=R_n(y_{n-1};\theta_n),\\
\text{变分法收敛} &\leftrightarrow& \Phi\text{ 的下降/最小化与第~\ref{thm:conditional-conv},\ref{thm:pgd-global} 条结论},\\
\text{离散统（一连续）} &\leftrightarrow& \sigma\text{（离散选代表）与 }\theta\text{（连续自由度）共同诱导的连续族}.
\end{array}
\]
\end{theorem}

\begin{Certificate}[证书：对象映射表与已证结论回填]\label{cert:alignment}
证书由两部分组成：
\begin{enumerate}
\item 对象级映射：调用定义~\ref{def:forget-map}--\ref{def:potential} 的符号对照；
\item 性质级映射：调用定理~\ref{thm:hole-iff-obtot}、\ref{thm:conditional-conv}、\ref{thm:pgd-global}。
\end{enumerate}
\end{Certificate}

\begin{proof}[证明（谓词因果链）]
由证书~\ref{cert:alignment}，逐项替换可得：
\[
\text{层化链每一环节}
\Rightarrow
\text{在本文对象集中有唯一对应接口}
\Rightarrow
\text{Chain-B 成为 Chain-A 的实例化}.
\]
定理成立。
\end{proof}

\subsection{“离散统（一连续）”的精确化}

\begin{lemma}[离散层阶与连续自由度的合成极限]\label{lem:discrete-continuous}
即使层阶 $n\in\mathbb N$ 为离散（可数），只要每层存在不可数参数 $\theta_n\in\Theta_n$，则由截面/选择规则构造出的代表族 $\sigma(x;\theta)$ 可在 $\theta$ 上形成连续族，
  从而出现“离散统（一连续）”现象。
\end{lemma}

\begin{Certificate}[证书：参数不可数性 + 选择规则正则性]\label{cert:discrete-continuous}
证书分两层：
\begin{enumerate}
\item 结构层：由定义~\ref{def:mixed-index}，参数方向不可数；
\item 正则层：附加“$\sigma$ 对参数连续”的正则性假设，则像族随参数连续变化。
\end{enumerate}
\end{Certificate}

\begin{proof}[证明（谓词因果链）]
由证书~\ref{cert:discrete-continuous}：
\[
|\Theta_n|\ge|\mathbb R|
\Rightarrow
\text{可用连续参数簇索引代表}
\Rightarrow
\sigma(x;\theta)
\text{可形成连续族（在连续性假设下）}.
\]
故“离散层阶 + 连续参数”可合成“离散统（一连续）”现象。
\end{proof}

\section{相对增益：从叙述链到可复用形式接口}

\begin{definition}[增益维度]\label{def:gain-metrics}
令 $\mathsf{OrigB}$ 表示仅含 Chain-B 的自然语言表达；$\mathsf{Align}$ 表示本文给出的对齐结构。
定义四维增益：
\[
\Delta_{\mathrm{ref}}:\text{可指称性},\quad
\Delta_{\mathrm{dis}}:\text{消歧},\quad
\Delta_{\mathrm{ver}}:\text{可验证性},\quad
\Delta_{\mathrm{imp}}:\text{可实现性}.
\]
\end{definition}

\begin{theorem}[相对 $\mathsf{OrigB}$ 的增益为正]\label{thm:gain-positive}
相对 $\mathsf{OrigB}$，本文结构 $\mathsf{Align}$ 在
\(
\Delta_{\mathrm{ref}},\Delta_{\mathrm{dis}},\Delta_{\mathrm{ver}},\Delta_{\mathrm{imp}}
\)
上均给出严格正增益。
\end{theorem}

\begin{Certificate}[证书：四维增益逐项见证]\label{cert:gain-positive}
证书按四维分解：
\begin{enumerate}
\item 见证对象：$\pi,\operatorname{Fib},\sigma,\Hole,\Ob,\Phi$ 的显式定义；
\item 见证语义：连续/离散、可数/不可数、崩溃/孔洞的一一对应；
\item 见证验证：收敛/稳态/全局最优被拆解为假设 + 结论；
\item 见证实现：诊断（孔洞）、修复（算子/约束）、优化（势函数）接口化。
\end{enumerate}
\end{Certificate}

\begin{proof}[证明（谓词因果链）]
由证书~\ref{cert:gain-positive}：
\[
\text{显式对象增加}
\Rightarrow
\Delta_{\mathrm{ref}}>0,
\quad
\text{语义一一映射增加}
\Rightarrow
\Delta_{\mathrm{dis}}>0,
\]
\[
\text{可证条款增加}
\Rightarrow
\Delta_{\mathrm{ver}}>0,
\quad
\text{接口落地增加}
\Rightarrow
\Delta_{\mathrm{imp}}>0.
\]
故四维增益均严格为正。定理成立。
\end{proof}

\begin{remark}
若把比较对象换成“已包含相同符号链与证明骨架的更强版本文本”，则增益量级会下降为“重包装/接口化增益”；本文不对该比较对象作额外假设。
\end{remark}

\section{结语：可审计的稳态构造模板}

本文给出一条可复用模板：

\begin{enumerate}
\item 用遗忘投影 $\pi$ 识别 $l_2$ 外延态的纤维冗余；
\item 用 $\mathcal A$ 与 $\Lift$ 定义可行域，并以 $\Hole$ 诊断崩溃点；
\item 若需要“孔洞 $\Leftrightarrow$ 上同调障碍”，引入塔式障碍与完备性假设；
\item 用 $\Ob$（可行性残差）与 $\Phi$（优化目标）分离建模；
\item 在明确的收缩/单调条件下推出条件化收敛；
\item 若要全局最优，再显式加入凸性/强凸性等强假设，否则标注为未证明。
\end{enumerate}

% ------------------------------
% 正文到此结束
% ------------------------------

\section*{参考文献}
\begin{thebibliography}{99}

\bibitem{Hatcher2002}
A.\ Hatcher,
\emph{Algebraic Topology},
Cambridge University Press, 2002.
\\ \url{https://pi.math.cornell.edu/~hatcher/AT/ATpage.html}

\bibitem{Whitehead1978}
G.\ W.\ Whitehead,
\emph{Elements of Homotopy Theory},
Springer, 1978.

\bibitem{Rudin1976}
W.\ Rudin,
\emph{Principles of Mathematical Analysis} (3rd ed.),
McGraw--Hill, 1976.

\bibitem{BoydVandenberghe2004}
S.\ Boyd and L.\ Vandenberghe,
\emph{Convex Optimization},
Cambridge University Press, 2004.
\\ \url{https://web.stanford.edu/~boyd/cvxbook/}

\bibitem{Nesterov2004}
Y.\ Nesterov,
\emph{Introductory Lectures on Convex Optimization: A Basic Course},
Kluwer Academic Publishers, 2004.

\bibitem{Beck2017}
A.\ Beck,
\emph{First-Order Methods in Optimization},
SIAM, 2017.

\bibitem{BauschkeCombettes2017}
H.\ H.\ Bauschke and P.\ L.\ Combettes,
\emph{Convex Analysis and Monotone Operator Theory in Hilbert Spaces} (2nd ed.),
Springer, 2017.

\bibitem{Dacorogna2008}
B.\ Dacorogna,
\emph{Direct Methods in the Calculus of Variations} (2nd ed.),
Springer, 2008.

\end{thebibliography}

\end{CJK*}
\end{document}
